% Einheit 1 von 4 – Ganzrationale Gleichungen (bank/gleichungen-loesen/e1.jsonl, zusammenbau v0.1)
%% TODO Zeitmarke Einheit 1: Marken-Zeile nicht lesbar
%% TODO Zweigzeile Teil 1: was hier gelernt wird (Satz in Schülersprache aus dem Einheitstitel) – Einheit 1
\einheitenkopf[e1]{Einheit 1 von 4 · Ganzrationale Gleichungen}
\zweigzeile{Abitur GK}
\setcounter{aufgabe}{10}

%% TODO Titel: Ich-kann-Satz fehlt in der Bank (Platzhalter: Kettenname) – Nr. 11
%% TODO Anweisung: ein Satz über der ersten Teilaufgabe fehlt in der Bank – Nr. 11
\begin{aufgabe}{Schnittpunkte und Stellen}
%% TODO \rechenplatz{n}: Schrittzahl des Grundfalls fehlt in der Bank (nur bei fester Schreibform, 2.3 b)
\begin{teile}
\teil Welcher Weg? $x^4 - 10x^2 + 9 = 0$ \kreuz{Lösungsformel} \\ \kreuz{Substitution} \\ \kreuz{Polynomdivision}
\end{teile}
\begin{gleichungsraster}[2]
\gl{\text{$f(x) = x^2 - 2x - 1$, $g(x) = x + 9$. Schnittpunkte?}} &
\gl{\text{Der Wasserstrahl eines Brunnens folgt $h(x) = -0{,}5x^2 + 2x + 1$ ($x$ waagerechter Abstand von der Düse, $h$ Höhe, beide in m). Er trifft im Fallen auf die Kante eines $1{,}5$ m hohen Beckens. Wie weit ist die Kante waagerecht von der Düse entfernt? Runde auf zwei Stellen. (FHR 2025)}} \\
\gl{\text{$f(x) = -0{,}5x^3 + 3x^2 - 2x + 6$. Der Punkt $P(0 | 6)$ liegt auf dem Graphen. An welchen weiteren Stellen hat $f$ den Wert 6? Runde auf zwei Stellen. (FHR 2021)}} &
\gl{\text{Der Querschnitt eines Bootsrumpfs wird oben von $g(x) = -x^2 + 1$ und unten von $f(x) = 0{,}25x^4 - x^2 - 3$ begrenzt. Koordinaten der Schnittpunkte? (FHR 2021)}} \\
\gl{\text{Ein Aufsteller wird oben von $g(x) = -0{,}25x^4 + 0{,}5x^2 + 3$ und unten von $f(x) = 0{,}25x^4 - 0{,}5x^2 - 1$ begrenzt. Koordinaten der Schnittpunkte? Kontrolle: eine Schnittstelle ist $x = 2$. (FHR 2022)}} &
\gl{\text{$f(x) = x^3 - 2x^2 - 3x + 4$, $g(x) = 2x - 2$. $Q(1 | 0)$ ist ein gemeinsamer Punkt. Koordinaten der beiden weiteren Schnittpunkte? (FHR 2022)}} \\
\gl{\text{Eine Spielhauswand hat oben den Bogen $f(x) = -0{,}25x^2 + 2x$ und darunter den Bogen $g(x) = 0{,}25x^2 - 2x + 6$ (1 LE = 1 m, Boden bei $y = 0$). Die Fläche unter $g$ zwischen den Schnittpunkten wird aus einer rechteckigen Holzplatte geschnitten, die bis zur Höhe der Schnittpunkte reicht. Schnittpunkte, Plattenmaße und Preis bei 40 € je m²? (FHR 2023)}} &
\end{gleichungsraster}
\begin{teile}
\teil $f(x) = x^2 + 3x - 1$, $g(x) = -x^2 - x + 5$. Zeige, dass sich die Graphen nur für $x = -3$ und $x = 1$ schneiden. \hfill (Abitur 2019 GK)
\teil Ein Blatt in einer Grafik wird von $u(x) = 0{,}5x^3$ und $v(x) = x^2 \cdot (3 - x)$ begrenzt. Zeige, dass $P(0 | 0)$ und $Q(2 | 4)$ die einzigen gemeinsamen Punkte sind. \hfill (Abitur 2020 GK)
\end{teile}
\begin{gleichungsraster}[2]
\gl{\text{$h(x) = \frac{4}{x - 2}$ ($x \ne 2$), $g(x) = x - 2$. Koordinaten der gemeinsamen Punkte? (Abitur 2019 GK)}} &
\end{gleichungsraster}
\begin{teile}
\teil Zeige: $-0{,}5 \cdot (x - 2) \cdot (x + 3)^2 = -0{,}5x^3 - 2x^2 + 1{,}5x + 9$. \hfill (Abitur 2020 GK)
\end{teile}
\begin{gleichungsraster}[2]
\gl{\text{Ein Waldstück wird von $f(x) = x^3 - x^2 - 6x + 3$ und $g(x) = x^2 - x - 3$ begrenzt; die Graphen schneiden sich in $Q(-2 | 3)$. Das Waldstück reicht von $Q$ bis zum nächsten Schnittpunkt $R$ rechts von $Q$; 1 LE entspricht 0,5 km. $x$-Koordinate von $R$ und Fläche des Waldstücks in km²? (FHR 2019)}} & \\
\end{gleichungsraster}
\end{aufgabe}

%% TODO Titel: Ich-kann-Satz fehlt in der Bank (Platzhalter: Kettenname) – Nr. 12
%% TODO Anweisung: ein Satz über der ersten Teilaufgabe fehlt in der Bank – Nr. 12
\begin{aufgabe}{Waagerechte Ausdehnung eines Profils in vorgegebener Höhe über die Stellen mit vorgegebenem Funktionswert berechnen}
\begin{gleichungsraster}[2]
\gl{\text{Der Längsschnitt eines Glases ist $f(x) = 0{,}05x^4$ für $-4 \le x \le 4$ (1 LE = 1 cm). 12 cm unter dem Rand verläuft eine Linie um das Glas. Wie breit ist das Glas dort, und wie lang ist die Linie? (Abitur 2017 LK)}} & \\
\end{gleichungsraster}
\end{aufgabe}

%% TODO Titel: Ich-kann-Satz fehlt in der Bank (Platzhalter: Kettenname) – Nr. 13
%% TODO Anweisung: ein Satz über der ersten Teilaufgabe fehlt in der Bank – Nr. 13
\begin{aufgabe}{Schnittpunkte und Stellen – Fehler finden, Begründen, Anwendung, Darstellungswechsel}
\begin{teile}
\teil Finde den Fehler und rechne richtig. \rechnung{x^4 - 7x^2 - 18 &= 0 &&\mid u = x^2 \\ u^2 - 7u - 18 &= 0 \\ u_1 = 9,\ u_2 &= -2 \\ x = \pm 3,\ x &= \pm\sqrt{2}}
\teil $x^2 + 3x - 10 = 0$ hat die Lösungen $2$ und $-5$. Begründe, warum es keine weitere Lösung gibt.
\end{teile}
\begin{gleichungsraster}[2]
\gl{\text{Eine Firma hat die Kosten $K(x) = x^2 + 2x + 16$ und den Erlös $E(x) = 12x$ ($x$ in hundert Stück, $K$ und $E$ in tausend €). Ab welcher und bis zu welcher Stückzahl macht sie keinen Verlust?}} &
\end{gleichungsraster}
\begin{teile}
\teil Die Abbildung zeigt die Graphen von $f(x) = (x - 1)^2 - 2$ und $g(x) = -x + 1$. Lies die Lösungen von $f(x) = g(x)$ und die Schnittpunkte ab.

\begin{ksys}[xmin=-3,xmax=4,ymin=-3,ymax=5,ablesen] \parabel{1}{1}{-2}{f} \gerade{-1}{1}{g} \end{ksys}
S1( \leerfeld | \leerfeld ), S2( \leerfeld | \leerfeld )
\end{teile}
\end{aufgabe}
